% Generated draft: requires mathematical and source review.
\documentclass[11pt]{article}
\usepackage[margin=25mm]{geometry}
\usepackage{fontspec}
\setmainfont{Latin Modern Roman}
\usepackage{amsmath,amssymb,mathtools,unicode-math}
\setmathfont{Latin Modern Math}
\usepackage{enumitem,needspace,xurl,hyperref}
\setlength{\parindent}{0pt}
\setlength{\parskip}{5pt}
\setlength{\emergencystretch}{3em}
\newcommand{\N}{\mathbb{N}}
\newcommand{\Z}{\mathbb{Z}}
\newcommand{\Q}{\mathbb{Q}}
\newcommand{\R}{\mathbb{R}}
\newcommand{\C}{\mathbb{C}}
\newcommand{\F}{\mathbb{F}}
\newcommand{\A}{\mathbb{A}}
\newcommand{\PP}{\mathbb{P}}
\newcommand{\E}{\mathbb{E}}
\newcommand{\eps}{\varepsilon}
\newcommand{\dd}{\,\mathrm{d}}
\begin{document}
\section*{UnsolvedMath ID: 30003086}
\textbf{Source problem code:} catalogue entry
\subsection{Definitions and mathematical statement}
Let \(d\) denote the number of lines in a configuration, and let \(\mathcal C_d(K)\) denote the class of configurations of \(d\) lines over a field \(K\) intended by the source. The supplied record does not specify the ambient space, whether the lines must be distinct, or the admissible range of \(d\). It also does not define the quantity \(H(K,L)\).

Subject to these unresolved definitions, the record specifies the \emph{absolute linear Harbourne constant} by
\[
H(d):=\min\bigl\{H(K,L): K\text{ is a field},\ L\in\mathcal C_d(K)\bigr\}.
\]
The quantification over fields has no stated restriction on characteristic. The use of a minimum is retained from the source; nonemptiness of the admissible class and attainment of the minimum are not established in the supplied record.

The requested problem is to determine \(H(d)\), or obtain estimates for it, as \(d\) varies over the intended admissible line counts. This is a conditional formulation: a fully specified mathematical problem requires the definition of \(H(K,L)\), the precise configuration class, and the admissible values of \(d\).
\subsection{Short English statement}
Determine or bound the smallest value of \(H(K,L)\) among all configurations of \(d\) lines over all fields. The supplied record omits the definition of \(H(K,L)\) and the precise admissibility conditions, so those details require source verification before the statement is mathematically complete.
\subsection{Sources}
\begin{itemize}
\item Local UnsolvedMath registry: supplied record with permanent ID \(30003086\), title \emph{Absolute Linear Harbourne Constants}, source problem code \texttt{catalogue entry}, in the Oberwolfach Reports Open Problems collection. The record gives proposed year 2016 and names no individual proposer.
\item Supplied report citation: \emph{Mini-Workshop: Arrangements of Subvarieties, and their Applications in Algebraic Geometry} (2016), Oberwolfach Reports, DOI 10.4171/owr/2016/14, pp.\ 653--702. \url{https://doi.org/10.4171/owr/2016/14}. The supplied background describes these reports as documenting mathematical workshops and their open-problem discussions.
\item Supplied literature reference, labeled ``Dumnicki--Harrer--Szpond'': \url{https://arxiv.org/abs/1507.04080}.
\item Supplied literature reference, labeled ``Dumnicki--Harrer--Szpond 2018'': \url{https://doi.org/10.1016/j.ffa.2018.03.001}.
\item Literature assessment dated exactly \texttt{2026-08-21}: the supplied curation labels the problem partially solved and states that values and bounds are known through 31 lines and for finite-projective-plane sizes, while the general value problem associated with an extracted conjectural passage remains unresolved. This is supplied dated curation derived from the OWR material, not an independently verified literature assessment; its status and cited claims require independent verification.
\item The registry labels the statement \texttt{corrected\_verified} and explains that notation, typography, and sentence structure were repaired without changing mathematical content. This is a supplied editorial label, not independent verification of the formulation or its completeness.
\end{itemize}
\end{document}
